\chapter{React Project Structure}

A flat \texttt{src/} folder works for a to-do demo. It falls apart the moment
your app has auth, protected routes, and more than three pages. Here is the
structure used throughout this book for the Task Management System.

\begin{diagrambox}[frontend/src Structure]
\begin{verbatim}
frontend/src/
|-- api/
|   `-- axios.js          (single configured axios instance)
|-- services/
|   |-- authService.js    (login, register, logout)
|   `-- taskService.js    (getTasks, createTask, updateTask...)
|-- components/
|   |-- TaskCard.jsx
|   |-- TaskList.jsx
|   `-- Navbar.jsx
|-- pages/
|   |-- Dashboard.jsx
|   |-- Login.jsx
|   `-- TaskDetails.jsx
|-- layouts/
|   `-- MainLayout.jsx    (Navbar + Sidebar + <Outlet/>)
|-- hooks/
|   `-- useAuth.js
|-- context/
|   `-- AuthContext.jsx
|-- utils/
|   `-- formatDate.js
|-- assets/
|   `-- logo.svg
|-- App.jsx
`-- main.jsx
\end{verbatim}
\end{diagrambox}

\section{What Goes Where}

\begin{table}[h]
\centering
\begin{tabularx}{\textwidth}{lXX}
\toprule
\textbf{Folder} & \textbf{Contains} & \textbf{Should NOT contain} \\
\midrule
\texttt{components/} & Small, reusable, presentation-focused UI pieces (\texttt{TaskCard}, \texttt{Button}, \texttt{Navbar}) & Route definitions, data fetching logic, page-level layout \\
\texttt{pages/} & Route-level components that compose components into a full screen (\texttt{Dashboard}, \texttt{Login}) & Low-level reusable UI meant to be shared across pages \\
\texttt{layouts/} & Shell components wrapping multiple pages (header/sidebar + \texttt{<Outlet/>}) & Business logic or data fetching for a specific feature \\
\texttt{hooks/} & Custom hooks encapsulating reusable stateful logic (\texttt{useAuth}, \texttt{useDebounce}) & JSX markup / rendered UI \\
\texttt{context/} & React Context providers and their state (\texttt{AuthContext}, \texttt{ThemeContext}) & One-off local component state that doesn't need to be global \\
\texttt{services/} & Functions that call the \texttt{api/} layer for a specific resource (\texttt{taskService.getTasks}) & Raw axios configuration, JSX \\
\texttt{api/} & The axios instance itself: base URL, interceptors, headers & Business logic specific to one resource \\
\texttt{utils/} & Pure helper functions with no React or network dependency (\texttt{formatDate}, \texttt{truncate}) & Components, hooks, anything using \texttt{useState} \\
\bottomrule
\end{tabularx}
\end{table}

\begin{notebox}[NOTE]
The distinction between \texttt{api/} and \texttt{services/} matters: \texttt{api/axios.js}
knows nothing about ``tasks'' or ``users'' -- it's generic HTTP plumbing.
\texttt{services/taskService.js} knows the specific endpoints and payload
shapes for tasks. This separation means changing your backend's base URL
never requires touching business logic files.
\end{notebox}

\begin{tipbox}[TIP]
A quick test for ``does this belong in \texttt{components/} or \texttt{pages/}?'':
if it needs its own route in \texttt{App.jsx}, it's a page. If it's used inside
two or more pages, it's a component.
\end{tipbox}
